\section{A Feasible Certificate and Its Finite-Precision Boundary}
\label{sec:certificates}

An approximate matrix output cannot be certified merely by reporting a polar
residual and a tangent residual.  If the tangent residual is nonzero, the
matrix is not feasible for the primal oracle, so its objective value need not
be a valid primal lower bound.

Assume here that \(\Theta=uv^\top\) is a unit rank-one normal.  Given an
arbitrary approximate matrix \(Z\), define the spectral normalization
\begin{equation}
\bar Z
=
\frac{Z}{\max\{1,\snorm{Z}\}},
\qquad
c=\ip{\Theta}{\bar Z},
\tag{54}
\label{eq:Zbar-c}
\end{equation}
and the feasible repair
\begin{equation}
\widetilde Z
=
\frac{\bar Z-c\Theta}{1+|c|}.
\tag{55}
\label{eq:Ztilde}
\end{equation}

\begin{theorem}[Feasible repair and valid primal--dual gap in exact arithmetic]
\label{thm:T6}
The repaired matrix satisfies
\begin{equation}
\ip{\Theta}{\widetilde Z}=0,
\qquad
\snorm{\widetilde Z}\le1.
\tag{56}
\label{eq:T6-feas}
\end{equation}
For any multiplier \(\lambda\), define
\begin{equation}
\Gap(\lambda,Z)
=
\nnorm{G+\lambda\Theta}
-\ip{G}{\widetilde Z}.
\tag{57}
\label{eq:true-gap}
\end{equation}
Then
\begin{equation}
\Gap(\lambda,Z)
\ge
p^\star-\ip{G}{\widetilde Z}
\ge0.
\tag{58}
\label{eq:gap-cover}
\end{equation}
Thus \eqref{eq:true-gap} is a valid upper bound on the primal LMO
suboptimality.

If
\begin{equation}
r_{\mathrm{pol}}
=
\nnorm{G+\lambda\Theta}
-\ip{G+\lambda\Theta}{\bar Z},
\tag{59}
\label{eq:rpol}
\end{equation}
then \(r_{\mathrm{pol}}\ge0\), and
\begin{equation}
\Gap(\lambda,Z)
\le
r_{\mathrm{pol}}
+|\lambda||c|
+\frac{2\nnorm{G}|c|}{1+|c|}.
\tag{60}
\label{eq:gap-decomposition}
\end{equation}
Finally, if a computable quantity \(\widehat\phi(\lambda)\) has been proved to
satisfy
\[
\widehat\phi(\lambda)\ge\nnorm{G+\lambda\Theta},
\]
then
\begin{equation}
\widehat\phi(\lambda)-\ip{G}{\widetilde Z}
\tag{61}
\label{eq:upper-certificate}
\end{equation}
is also a valid certificate.
\end{theorem}

\begin{proofidea}
First normalize into the spectral ball.  Then subtract the normal component
and divide by a worst-case triangle-inequality factor so that tangency and
spectral feasibility both hold exactly.  Only after this repair is the
dual-minus-primal difference a legitimate gap.
\end{proofidea}

\begin{intuition}
The repair turns any approximate polar map---including one affected by truncation,
finite precision, or an imperfect multiplier---into a direction that the
inner oracle may accept at a prescribed objective accuracy.  The certificate measures the price of the
repair rather than silently treating an infeasible matrix as primal feasible.
\end{intuition}

\begin{proof}
By definition, \(\snorm{\bar Z}\le1\).  Because the rank-one unit normal
satisfies \(\ip{\Theta}{\Theta}=\fnorm{\Theta}^2=1\),
\[
\ip{\Theta}{\widetilde Z}
=
\frac{c-c}{1+|c|}
=0.
\]
Furthermore,
\[
\snorm{\bar Z-c\Theta}
\le
\snorm{\bar Z}+|c|\snorm{\Theta}
\le1+|c|.
\]
Dividing by \(1+|c|\) proves \eqref{eq:T6-feas}.

Weak duality from Theorem~\ref{thm:T1} gives
\[
\nnorm{G+\lambda\Theta}\ge p^\star.
\]
Since \(\widetilde Z\) is primal feasible,
\[
\ip{G}{\widetilde Z}\le p^\star.
\]
Subtracting proves \eqref{eq:gap-cover}.

For the decomposition, put \(A=G+\lambda\Theta\).  Adding and subtracting
\(\ip{A}{\bar Z}\) gives
\[
\begin{aligned}
\Gap(\lambda,Z)
&=
\nnorm{A}-\ip{G}{\widetilde Z}\\
&=
r_{\mathrm{pol}}
+\lambda c
+\ip{G}{\bar Z-\widetilde Z}.
\end{aligned}
\]
Equation \eqref{eq:Ztilde} gives
\[
\bar Z-\widetilde Z
=
\frac{|c|\bar Z+c\Theta}{1+|c|}.
\]
By spectral/nuclear duality,
\[
\begin{aligned}
\left|\ip{G}{\bar Z-\widetilde Z}\right|
&\le
\frac{|c|\nnorm{G}\snorm{\bar Z}
      +|c|\nnorm{G}\snorm{\Theta}}
     {1+|c|}\\
&\le
\frac{2\nnorm{G}|c|}{1+|c|}.
\end{aligned}
\]
Using \(\lambda c\le|\lambda||c|\) proves
\eqref{eq:gap-decomposition}.  Replacing the exact dual value by any proved
upper bound preserves the coverage inequality, giving
\eqref{eq:upper-certificate}.
\end{proof}

\begin{remark}[The same repair for an arbitrary nonzero normal]
\label{rem:general-normal-repair}
The rank-one normalization makes the formulas particularly simple but is
not essential.  For any nonzero \(\Theta\), normalize \(Z\) as in
\eqref{eq:Zbar-c} and put
\[
 b=\ip{\Theta}{\bar Z},\qquad
 d=\frac{b}{\|\Theta\|_F^2},\qquad
 \eta=|d|\|\Theta\|_2,\qquad
 \widetilde Z=\frac{\bar Z-d\Theta}{1+\eta}.
\]
Then \(\langle\Theta,\widetilde Z\rangle=0\) and
\(\|\widetilde Z\|_2\le1\).  With
\(r_{\rm pol}=\|G+\lambda\Theta\|_*
-\langle G+\lambda\Theta,\bar Z\rangle\), the same calculation gives
\[
 0\le p^\star-\langle G,\widetilde Z\rangle
 \le\|G+\lambda\Theta\|_*-\langle G,\widetilde Z\rangle
 \le r_{\rm pol}+|\lambda||b|
       +\frac{2\|G\|_*\eta}{1+\eta}.
\]
Indeed, \(\bar Z-\widetilde Z=(\eta\bar Z+d\Theta)/(1+\eta)\), whose
spectral norm is at most \(2\eta/(1+\eta)\).  This is an elementary
extension of the same feasibility argument, not a further smoothing law.
\end{remark}

\begin{boundary}
The feasibility and coverage statements above are exact-real-arithmetic
identities.  A machine-rigorous certificate additionally needs outward-rounded
upper bounds for the nuclear norm and spectral normalization, a downward-
rounded primal objective, and enclosures for the tangency calculation.  Plain
double-precision evaluation is a numerical check, not by itself a validated
certificate.  These are necessary ingredients, not a complete validated
implementation: an interval enclosure of the tangency residual alone does
not certify exact membership in the hyperplane unless the final repair is
itself performed exactly or with an interval-safe correction.
\end{boundary}

\begin{corollary}[Optional unit-spectral-norm output]
\label{cor:unit-repair}
Let
\[
P=\bar Z-c\Theta.
\]
If \(P\ne0\), choose
\begin{equation}
s_G=
\begin{cases}
1,&\ip{G}{P}\ge0,\\
-1,&\ip{G}{P}<0,
\end{cases}
\qquad
Z_{\rm unit}=s_G\frac{P}{\snorm{P}}.
\tag{62}
\label{eq:unit-repair}
\end{equation}
Then
\[
\ip{\Theta}{Z_{\rm unit}}=0,\qquad
\snorm{Z_{\rm unit}}=1,
\]
and
\begin{equation}
\ip{G}{Z_{\rm unit}}\ge\ip{G}{\widetilde Z}.
\tag{63}
\label{eq:unit-repair-objective}
\end{equation}
Consequently, replacing \(\widetilde Z\) by \(Z_{\rm unit}\) never worsens the
valid primal--dual certificate.

Likewise, if \(p^\star>0\), then the smooth solution
\(\Phi_\delta\ne0\) may be replaced by
\[
\Phi_\delta^{\rm unit}
=\frac{\Phi_\delta}{\snorm{\Phi_\delta}}.
\]
It remains tangent, has unit spectral norm, and has no smaller original
linear objective.  The lower bound
\[
\sigma_1(G+\lambda_\delta\Theta)\ge s_0:=p^\star/q>0,
\qquad q=\min\{m,n\},
\]
holds uniformly for every \(\delta>0\), because
\(q\sigma_1(G+\lambda\Theta)\ge\|G+\lambda\Theta\|_*\ge p^\star\)
for every \(\lambda\).  Therefore the rescaling factor satisfies
\begin{equation}
\frac{1}{\snorm{\Phi_\delta}}
=
\sqrt{1+\frac{\delta^2}
{\sigma_1(G+\lambda_\delta\Theta)^2}}
=1+O(\delta^2).
\tag{64}
\label{eq:smooth-unit-factor}
\end{equation}
\end{corollary}

\begin{proof}
Tangency and unit norm in \eqref{eq:unit-repair} are immediate.  Moreover,
\(\snorm{P}\le1+|c|\), so if \(\ip{G}{P}\ge0\),
\[
\ip{G}{Z_{\rm unit}}
=\frac{\ip{G}{P}}{\snorm{P}}
\ge\frac{\ip{G}{P}}{1+|c|}
=\ip{G}{\widetilde Z}.
\]
If \(\ip{G}{P}<0\), the left side is positive whereas the right side is
negative.

For the smooth direction, comparison with the feasible zero direction in
\eqref{eq:regularized-primal} gives
\[
\ip{G}{\Phi_\delta}
\ge\delta R(\Phi_\delta)\ge0.
\]
To see rigorously that \(\Phi_\delta\ne0\) when \(p^\star>0\), let
\(\Phi^\star\) be an exact primal optimizer.  For sufficiently small
\(t>0\), the exact identity
\(1-\sqrt{1-z^2}=z^2/(1+\sqrt{1-z^2})\le z^2\) gives
\[
 R(t\Phi^\star)\le t^2\|\Phi^\star\|_F^2.
\]
Consequently
\[
 \ip{G}{t\Phi^\star}-\delta R(t\Phi^\star)
 \ge tp^\star-\delta t^2\|\Phi^\star\|_F^2>0,
\]
whereas the regularized objective at zero is zero.  Hence the unique smooth
optimizer is nonzero, and positive radial rescaling cannot decrease its
linear objective.  Finally,
\eqref{eq:Zdelta} gives
\[
\snorm{\Phi_\delta}
=
\frac{\sigma_1(G+\lambda_\delta\Theta)}
{\sqrt{\sigma_1(G+\lambda_\delta\Theta)^2+\delta^2}},
\]
which proves \eqref{eq:smooth-unit-factor}.
\end{proof}

\begin{boundary}
The unit repair is certified primal post-processing and an experimental
scale-control device.  The map \(P\mapsto P/\snorm{P}\) is discontinuous at
\(P=0\) and may amplify numerical noise nearby.  Moreover,
\(\Phi_\delta^{\rm unit}\) is no longer the exact optimizer of the
regularized primal at finite \(\delta\), so its Fenchel center-path identity
must not be asserted pointwise.  For a fixed instance with \(p^\star>0\),
however, its rescaling is \(1+O(\delta^2)\)
and it has the same limiting center.
\end{boundary}

\subsection{How to read the decomposition}

The three terms in \eqref{eq:gap-decomposition} have different meanings.

\begin{enumerate}
  \item \(r_{\mathrm{pol}}\) measures matrix-function or polar optimality
        error.
  \item \(|\lambda||c|\) measures the coupling between multiplier size and
        tangent residual.
  \item The last term is the explicit price of moving the infeasible output
        into the primal feasible set.
\end{enumerate}

\paragraph{Certified stopping rule.}
Given a target inner error budget \(\varepsilon_t\), stop the inner oracle only
when a valid quantity of the form \eqref{eq:true-gap} or
\eqref{eq:upper-certificate} is at most \(\varepsilon_t\).  Reporting only
\(|c|\) is not sufficient.

The matrix \(Z\) in these diagnostics is the actual returned output of the
implemented map.  Replacing it by \(\Sc(G+\lambda\Theta)\), or computing
the tangent residual with an estimated normal while interpreting it as a
residual for the true normal, changes the quantity being assessed.  A small
tangent residual alone also does not imply a small \(r_{\rm pol}\) or a
small feasible gap.  Conversely, a small computed floating-point gap is
evidence for numerical accuracy only within its recorded norm and pairing
errors; it is not a machine-rigorous certificate without validated bounds.

\subsection{A certifiable reference homotopy oracle}

The theory suggests the following implementation pattern.

\begin{enumerate}
  \item \textbf{Choose a smoothing level.}
        Start from a moderate \(\delta\); warm-start it from the previous
        oracle call when appropriate.
  \item \textbf{Bracket the smooth multiplier.}
        Use the proved multiplier bound
        \eqref{eq:smooth-lambda-bound}, or a tighter bracket whose coverage has
        been verified.
  \item \textbf{Solve the continuous root.}
        Apply bisection or a safeguarded scalar solver to
        \(h_\delta(\lambda)=0\).  Unlike the compact selection, this function
        is continuous and strictly increasing.
  \item \textbf{Form the smooth direction.}
        Evaluate \(Z_\delta(G+\lambda\Theta)\).  If the scalar root is only
        approximate or the matrix function is computed in finite precision,
        apply \eqref{eq:Zbar-c}--\eqref{eq:Ztilde}.
  \item \textbf{Certify.}
        Compute \eqref{eq:true-gap}, or a proved upper version
        \eqref{eq:upper-certificate}.
  \item \textbf{Refine only when needed.}
        If the certificate exceeds the requested tolerance \(\varepsilon_t\), reduce
        \(\delta\), tighten the scalar root, or switch to the exact
        set-valued SVD oracle.  Otherwise accept the repaired direction.
\end{enumerate}

This homotopy is ``accuracy on demand'' only after the normal accuracy and any
assumptions required by a downstream method have been certified independently.
\Cref{thm:T6} makes only the \emph{inner} oracle gap computable; it does not certify the
normal or an outer descent model.
Consequently this is a certifiable reference workflow, not a claim of a
fully deployable per-step large-model routine.

\begin{remark}[Conservatism]
The repair is deliberately simple and globally safe.  It may be conservative:
projecting onto the exact intersection of a spectral ball and a hyperplane
could produce a tighter primal value, but would itself require a nontrivial
optimization.  The displayed repair prioritizes a closed-form guarantee.
\end{remark}
